\CUSTOMappendix{Система сбора истории камунды}

\begin{lstlisting}[
  caption={Контроллер},
  label={lst:activity-log}
]
@Slf4j
@RestController
@RequestMapping("/api")
@RequiredArgsConstructor
@Tag(name = "Activity Log API")
public class ActivityLogController {

  private final ActivityNameRegistry activityNameRegistry;

  /**
   * Скачивание журнала активностей в формате Excel.
   *
   * @param startTime Время начала периода (необязательно). Формат: yyyy-MM-dd'T'HH:mm:ss
   * @param endTime Время окончания периода (необязательно). Формат: yyyy-MM-dd'T'HH:mm:ss
   * @return Excel-файл с журналом активностей
   */
  @GetMapping("/activity-log")
  @Operation(
      summary = "Скачать журнал активностей",
      description =
          "Возвращает Excel-файл с журналом активностей Camunda за указанный период. "
              + "Если параметры времени не указаны, возвращается полный журнал.")
  @ApiResponses(
      value = {
        @ApiResponse(
            responseCode = "200",
            description = "Файл успешно сгенерирован и возвращен",
            content =
                @Content(
                    mediaType =
                        "application/vnd.openxmlformats-officedocument.spreadsheetml.sheet")),
        @ApiResponse(responseCode = "400", description = "Некорректные параметры запроса"),
        @ApiResponse(
            responseCode = "404",
            description = "Нет данных для отображения в журнале активностей за указанный период"),
        @ApiResponse(responseCode = "500", description = "Внутренняя ошибка сервера")
      })
  public ResponseEntity<Resource> downloadActivityLog(
      @RequestParam(required = false) @DateTimeFormat(pattern = "yyyy-MM-dd'T'HH:mm:ss")
          Date startTime,
      @RequestParam(required = false) @DateTimeFormat(pattern = "yyyy-MM-dd'T'HH:mm:ss")
          Date endTime) {

    Date currentTime = new Date();

    if ((startTime != null && endTime != null && startTime.after(endTime))
        || (startTime != null && startTime.after(currentTime))
        || (endTime != null && endTime.after(currentTime))) {
      return ResponseEntity.badRequest().build();
    }
    try {

      log.info("Запрос на получение логгов по камунде");
      File file = activityNameRegistry.generateActivityLogFile(startTime, endTime);

      if (file == null) {
        log.warn("Нет данных для отображения в журнале активностей");
        return ResponseEntity.notFound().build();
      }

      InputStreamResource resource = new InputStreamResource(new FileInputStream(file));

      DateFormat dateFormat = new SimpleDateFormat("yyyy-MM-dd_HH-mm-ss");
      String formattedDate = dateFormat.format(new Date());

      HttpHeaders headers = new HttpHeaders();
      headers.add(
          HttpHeaders.CONTENT_DISPOSITION,
          "attachment; filename=ActivityLog_" + formattedDate + ".xlsx");

      return ResponseEntity.ok()
          .headers(headers)
          .contentLength(file.length())
          .contentType(
              MediaType.parseMediaType(
                  "application/vnd.openxmlformats-officedocument.spreadsheetml.sheet"))
          .body(resource);
    } catch (IOException e) {
      log.error("Ошибка при создании файла журнала активностей", e);
      return ResponseEntity.internalServerError().build();
    }
  }
}
\end{lstlisting}


\begin{lstlisting}[
  caption={Сервис},
  label={lst:activity-log}
]
/**
 * ActivityNameRegistry загружает отображения ID активностей в человекочитаемые имена из JSON-файла
 * при старте приложения.
 */
@Component
@RequiredArgsConstructor
public class ActivityNameRegistry {
  private static final Logger LOGGER = LoggerFactory.getLogger(ActivityNameRegistry.class);

  private Map<String, ActivityInfo> activityInfoMap = new HashMap<>();

  private final HistoricActivityInstanceApi historicActivityRestService;

  @Value("${activity.decor.path}")
  private String pathToActivityDecor;

  /** Инициализирует отображения из JSON-файла при старте Spring-контейнера. */
  @PostConstruct
  public void init() {
    try {
      ClassPathResource resource = new ClassPathResource(pathToActivityDecor);
      ObjectMapper mapper = new ObjectMapper();
      this.activityInfoMap = mapper.readValue(resource.getInputStream(),
          new TypeReference<Map<String, ActivityInfo>>() {});
      LOGGER.info("Loaded {} activity info mappings");
    } catch (IOException e) {
      LOGGER.error("Failed to load activity-decor.json", e);
    }
  }

  /** Возвращает человекочитаемое имя активности по её ID. */
  public String getName(String activityId) {
    ActivityInfo info = this.activityInfoMap.get(activityId);
    return info != null ? info.getName() : activityId;
  }

  /** Возвращает схему активности по её ID. */
  public String getSchema(String activityId) {
    ActivityInfo info = this.activityInfoMap.get(activityId);
    return info != null ? info.getSchema() : null;
  }

  /** Класс для хранения информации об активности. */
  @Getter
  @Setter
  public static class ActivityInfo {
    private String name;
    private String schema;
  }

  /**
   * Генерирует активити лог.
   *
   * @param startTime Optional start date filter
   * @param endTime Optional end date filter
   * @return Generated Excel file with activity log
   * @throws IOException If file cannot be created
   */
  public File generateActivityLogFile(final Date startTime, final Date endTime)
      throws IOException {
    HistoricActivityInstanceQueryDto query = new HistoricActivityInstanceQueryDto();
    Optional.ofNullable(startTime).ifPresent(query::setStartedAfter);
    Optional.ofNullable(endTime).ifPresent(query::setStartedBefore);

    List<HistoricActivityInstanceDto> activities = new ArrayList<>();

    try {
      activities = historicActivityRestService
          .queryHistoricActivityInstances(null, null, query).stream()
          .filter(h -> h.getEndTime() != null)
          .collect(Collectors.toList());

      if (activities.isEmpty()) {
        return null;
      }

      Workbook workbook = new XSSFWorkbook();
      Sheet sheet = workbook.createSheet("Activity Log");
      int rowNum = 1;

      Row headerRow = sheet.createRow(0);
      headerRow.createCell(0).setCellValue("Activity ID");
      headerRow.createCell(1).setCellValue("Activity Name");
      headerRow.createCell(2).setCellValue("Start Time");
      headerRow.createCell(3).setCellValue("End Time");
      headerRow.createCell(4).setCellValue("Duration (ms)");
      headerRow.createCell(5).setCellValue("Process Instance ID");
      headerRow.createCell(6).setCellValue("Schema");
      headerRow.createCell(7).setCellValue("Parent Activity Instance Id");
      headerRow.createCell(8).setCellValue("Parent Activity Instance Name");
      headerRow.createCell(9).setCellValue("Called Process Instance Id");
      headerRow.createCell(10).setCellValue("Root Process Instance Id");
      headerRow.createCell(11).setCellValue("Parent Activity Instance Schema");
      headerRow.createCell(12).setCellValue("Activity Type");


      SimpleDateFormat timeFormat = new SimpleDateFormat("yyyy-MM-dd HH:mm:ss.SSS");

      for (HistoricActivityInstanceDto activity : activities) {
        Row row = sheet.createRow(rowNum++);
        row.createCell(0).setCellValue(activity.getActivityId());
        row.createCell(1).setCellValue(getName(activity.getActivityId()));
        row.createCell(2).setCellValue(timeFormat.format(activity.getStartTime()));
        row.createCell(3).setCellValue(timeFormat.format(activity.getEndTime()));
        row.createCell(4).setCellValue(
            activity.getDurationInMillis() != null ? activity.getDurationInMillis() : 0);
        row.createCell(5).setCellValue(activity.getProcessInstanceId());
        row.createCell(6).setCellValue(getSchema(activity.getActivityId()));
        row.createCell(7).setCellValue(activity.getParentActivityInstanceId());
        row.createCell(8).setCellValue(getName(
            activity.getParentActivityInstanceId()));
        row.createCell(9).setCellValue(activity.getCalledProcessInstanceId());
        row.createCell(10).setCellValue(activity.getRootProcessInstanceId());
        row.createCell(11).setCellValue(getSchema(
            activity.getParentActivityInstanceId()));
        row.createCell(12).setCellValue(activity.getActivityType());
      }

      File tempFile = File.createTempFile("activityLog", ".xlsx");
      try (FileOutputStream outputStream = new FileOutputStream(tempFile)) {
        workbook.write(outputStream);
      }

      workbook.close();
      return tempFile;

    } catch (ApiException e) {
      LOGGER.error(e.getResponseBody(), e);
      throw new IOException("Failed to retrieve activity data", e);
    }
  }
}
\end{lstlisting}
